\documentclass[letterpaper,10pt]{texMemo}

\usepackage{graphicx}

\memoto{Dr. Mark Yoder}
\memofrom{Michael McDonald}
\memosubject{Week 5 (Kernel and memmap, Exercises 8, 8a, 11b, 23) Memo}
\memodate{\today}
\logo{\includegraphics[width=0.5\textwidth]{RoseLogo.png}}

\begin{document}
\maketitle

\section{Exercise Summary}

\noindent Homework 5 is complete.

\begin{center}
\begin{tabular}{|p{4cm}|l|l|p{8cm}|}
\hline
\textbf{Objective} & \textbf{Exercise} &\textbf{Status} & \textbf{Notes}\\
\hline
Compile the 3.8 Kernel & Exercise 8 & Completed & I did this in Memo 2,
because it was on the calendar then.\\
\hline
Install necessary cross compilation tools & Exercise 8a & Completed & I did this
in Memo 2,
because it was on the calendar then.\\
\hline
Memmap GPIO access & Exercise 11 & Completed & This is awesome. It also makes
me want to write another layer of abstraction to do GPIO calls, because it's a
bit of a pain.\\
\hline
Building the Linux Kernel & Exercise 23 & Completed & Can't wait to start
playing with the Kernel and changing it :) \\
\hline
\end{tabular}
\end{center}

\noindent I ran the experiments in the extras section and experienced some
strange behavior. First off, the performance seemed to actually be on par with
the shell file (if not worse). The output signals weren't clean, and they were
definitely not accurately representing the input properly. They also delayed
significantly longer than anticipated, from nearly zero to two or three ms,
which is a significant delay that we should not be dealing with. Setting the
nice priority (\verb nice  \verb -n  \verb .\gpioThru  \verb & ) didn't seem to
change things that much either, even when I set the nice to -20
(on par with the Kernel level processes). I wasn't running any other processes
in my Userspace, so I'm not entirely sure what is causing my problems, but I
would like to test it on another bone to see if it's my bone or if it's my
code.\\


\noindent The main points I took away from these exercises were the following:
\begin{itemize}
 \item I'm really fascinated at how the Linux kernel was created, and how it
accomplishes the tasks it needs to. I'd also be curious to know how small one
could make it in order to put it on smaller and smaller platforms (possibly
even cut into the market of Arduino [especially now that those are starting to
grow in size]).
  \item I feel like taking this class, especially checking out all the
low level code, makes computing feel less high level and abstract, and adds a
concreteness that 
\end{itemize}

\end{document}
